% Physical page 5 onward: the developed commentary of the concise study,
% organized by the questions the three interpretations answer differently and
% not by the order of the Mass. Each unit draws on several appointed elements
% and sets the witnesses' answers beside one another.

\section{The Propers: Detailed Commentary}
\zlabel{triptych:concise:commentary:start}

\subsection{What the king prepares, and for whom}

The parable's first sentence holds the first interpretation's question: who
is the king, who is the son, and what is the wedding? Gregory the Great
answers twice and corrects himself between. The Father made the wedding for
his Son, he says first, when in the Virgin's womb he joined him to human
nature; but such a joining is ordinarily made of two persons, and he will not
have the Redeemer thought of as made out of two. So he restates it:
\latin{Apertius ergo atque securius dici potest quia in hoc Pater regi Filio
nuptias fecit, quo ei per incarnationis mysterium sanctam Ecclesiam
sociavit}, it can be said more openly and more safely that the Father made the
wedding for the King his Son in this, that through the mystery of the
Incarnation he joined the holy Church to him (\work{Homiliae in Evangelia}
38.3).

Jerome says who that Church is: God almighty makes the wedding \latin{Domino
nostro Iesu Christo et Ecclesiae, quae tam ex Iudaeis, quam ex gentibus,
congregata est}, for our Lord Jesus Christ and the Church, gathered as much
from the Jews as from the gentiles (\work{Commentarii in Matthaeum} III, PL 26,
col.~159). Irenaeus, writing against those who divided the God who calls from
the God who judges, sees ``one King and Lord, the Father of all'', who ``had
from the beginning prepared the marriage for His Son'' and ``called together
from all the highways, that is, from all nations, [guests] to the marriage
feast of His Son'' (\work{Against Heresies} IV.36.5--6). John Chrysostom asks
why the parable speaks of a marriage at all, and answers, ``That thou mightest
learn God's tender care, His yearning towards us'', for there ``all things are
full of spiritual joy''; and because the parable comes after the vineyard and
its slain son, he adds of a wedding made after a death, ``Hereby He proclaimed
the resurrection also'' (\work{Homilies on Matthew} 69.1). Hilary of Poitiers
draws that thread tightest: the wedding is \latin{vitae coelestis et in
resurrectione suscipiendae aeternae gloriae sacramentum}, the sacrament of the
heavenly life and of the eternal glory to be received in the resurrection
(\work{Commentarius in Matthaeum} 22.3, PL 9, col.~1042). They agree that the
king is the Father, the son Christ, and that the call goes beyond the first
invited; they name the bride differently.

Isaiah's feast is the other half of the \work{Ordo}'s correlation. Jerome
places it after the Passion, for all the nations on Mount Sion, and reads the
wiping of tears as the moment \latin{quando morte superata, Christi advenerit
regnum}, when death is overcome and Christ's kingdom has come
(\work{Commentarii in Isaiam} VIII, PL 24, col.~300). Thomas Aquinas makes the
one join between the two readings that any witness cited makes. The feast is
made on this mountain \latin{quia ibi Christus passus est \ldots\ vel quia ibi
iudicabit, ut dicitur Matth. 22: Ecce prandium meum paravi}, because Christ
suffered there, or because he will judge there, as Matthew 22 says, ``Behold, I
have prepared my dinner''. He distinguishes three feasts: the household feast
of the Church on pilgrimage, for which he cites 1~Corinthians 11:26 and this
Sunday's psalm, \latin{Impinguasti in oleo caput meum}; the private feast of
the soul; and the solemn feast of the heavenly court, given \latin{ad
satietatem}, to the full (\work{Expositio super Isaiam} c.~25,
pp.~501--502). On the psalm he reads the table also as
\latin{mensam sacramentalem, scilicet altaris}, the sacramental table of the
altar (\work{Postilla super Psalmos} 22). That the feast is the altar's is his
reception; the literal sense of Isaiah and of the psalm is a banquet God
gives.

The other interpretations take the feast as given. The second asks how a guest
is fit for it. The third asks what the invited refused, and Gregory answers
that it was a mystery: the man given over to earthly labour or worldly
business \latin{mysterium incarnationis dominicae pensare et secundum illud
vivere dissimulat}, pretends not to see the need to weigh the mystery of the
Lord's Incarnation and to live by it (38.5). His restated wedding thus stands
at the centre of the first interpretation and of the third.

\subsection{Now, or at the resurrection?}

Where does the feast stand? Gregory places the parable's wedding in the
present Church, because of its mixed guests: \latin{per has regis nuptias
praesens Ecclesia designatur, in qua cum bonis et mali conveniunt}, the present
Church, in which the bad come together with the good (38.7); Luke's great
supper, from which no one is cast out, is for him another parable and the
eternal banquet (38.1). Augustine speaks of ``two feasts of the Lord; one to
which the good and evil come, the other to which the evil come not'', and joins
them: ``Be that feast which now is, received worthily, that we may attain to
the other'' (\work{Sermo} 90.1, 4). Hilary makes the wedding the glory of the
resurrection. Aquinas gives four expositions, among them the wedding
\latin{in communi resurrectione}, when what is mortal in us will be swallowed
up by life, but for the parable follows Gregory: \latin{secundum Gregorium,
oportet exponere de praesentibus}, it must be expounded of present things
(\work{Super Matthaeum} c.~22, p.~280).

These are emphases, not contradictions: Gregory does not read the parable's
wedding as the last banquet, nor Hilary as the Church on earth. The first
interpretation joins the two feasts without identifying them. In both forms
of the Gospel, what joins them is the resurrection, which Hilary's wedding awaits and Chrysostom hears in a wedding made after a
death. Isaiah supplies the end of the line: Jerome reads its death cast down
with Paul's ``Death is swallowed up in victory'' (1~Cor 15:54), and Aquinas
reads the wedding of the resurrection in Paul's words on the mortal swallowed
up by life (2~Cor 5:4). Where the longer form is read, the judgment stands
between as well, for every witness who comments on the king's coming in to see
the guests places it there, and the present feast becomes the anteroom of the
last.

\subsection{Why the invited would not come}

Chrysostom puts the third interpretation's question aloud: ``For who would not
choose to come to a marriage, and that a King's marriage, and of a King making
a marriage for a Son?'' (69.1). The first refusal is not hostility. The
invited ``neglected and went their ways, one to his farm and another to his
merchandise'' (22:5); only ``the rest'' kill the servants (22:6).

Two witnesses judge the excuses and two read them. Chrysostom, who lists
Luke's oxen and wife beside Matthew's farm, grants that ``the excuses seem to
be reasonable; but hence we learn, though the things which hinder us be
necessary, to set the things spiritual at a higher price than all'' (69.1).
Aquinas: \latin{Videbantur habere iustam causam exterius; sed Dominus non
recipit: quia nulla temporalia debent detinere de veniendo ad Deum}, outwardly
they seemed to have a just cause, but the Lord does not accept it, because no
temporal things ought to keep anyone from coming to God (p.~281). Gregory
reads the two excuses as two habits of life: to go to the farm is
\latin{labori terreno immoderate incumbere}, to lie on earthly labour beyond
measure, and to go to the merchandise \latin{actionum saecularium lucris
inhiare}, to gape after the profits of worldly business (38.5). Hilary reads
ambition and avarice, some held \latin{ambitione saeculi tamquam agro}, by
worldly ambition as by a field, more by trading \latin{ob pecuniae
cupiditatem}, for love of money (22.5, col.~1043). Gregory and Hilary name
different appetites; Chrysostom and Aquinas agree that even a reasonable one
is no excuse.

Jerome weighs the busy against the violent: those taken up with other things
are guilty of a lesser crime than those who \latin{verterunt humanitatem in
crudelitatem}, turned kindness into cruelty (col.~159). The lesser crime is
still a refusal; but his distinction marks this interpretation's limit, for a
story of distraction leaves out the parable's violent half, which the first
interpretation has to carry. The second command sends the servants \latin{ad
exitus viarum}. Jerome reads the outlets of the roads as the place of the
gentiles (col.~160); Gregory as the failure of men's actions, \latin{quia illi
plerumque facile ad Deum veniunt, quos in terrenis actibus prospera nulla
comitantur}, because those very often come easily to God whom no prosperity
accompanies in their earthly doings (38.6). Jerome answers who the new guests
were, Gregory why they came.

\subsection{The burnt city, and to whom the parable is spoken}

``The rest'' kill the servants, and the king's armies destroy them and burn
their city (22:6--7). Chrysostom reads the burning as ``the things that took
place under Vespasian and Titus'', forty years after the Passion, ``that He
might show His long suffering'' (69.1). Jerome gives two readings of the
armies, avenging angels or the Romans, and Aquinas gives both, adding that
God's anger \latin{non commotionem significat, sed vindictam}, signifies not
agitation but just punishment (p.~281). Gregory and Hilary read the armies as
angels and the fire as the eternal punishment of the persecutors (38.5). This
is the verse the first interpretation, which rejoices in tears wiped from
every face, must carry.

Irenaeus, Hilary, Chrysostom, Jerome and Aquinas identify the first invited
with Israel; Chrysostom says the parable foretells ``the casting out of the
Jews, and the calling of the Gentiles'', and Jerome's sentence on Isaiah sets
the feast \latin{nequaquam populo Iudaeorum}, not for the people of the Jews.
Gregory identifies the refusers with the persecutors of the preachers. The
parable was spoken to the chief priests and elders of the people (21:23). The
Second Vatican Council, in \work{Nostra aetate} 4, governs
how such readings are presented, though it passes no judgment on any Father's
exegesis. It recalls that the gifts and the calling of God are without
repentance (Rom 11:29), and teaches that what happened in the Passion cannot
be charged ``against all the Jews, without distinction, then alive, nor
against the Jews of today'' and that the Jews ``should not be presented as
rejected or accursed by God, as if this followed from the Holy Scriptures''.
Under that bound the identification belongs to the Fathers who make it, the
parable's hearers are the leaders in the temple, and the warning falls, as
Gregory and Chrysostom themselves turn it, on the baptized: \latin{Nos sumus,
fratres charissimi, qui in nuptiis Verbi discumbimus}, we are the ones, dearest
brethren, who recline at the wedding of the Word (38.11); ``Hear ye, as many
as having partaken of the mysteries'' (69.2). Isaiah's ``all people'' cannot be
narrowed to ``not this people'' on the strength of Jerome's phrase, which is
his alone.

\subsection{The guest without a wedding garment}

Where the longer form is read, the second interpretation meets the parable's
sharpest difficulty: the king's own servants gathered the guest, so why is he
blamed for his clothes? The witnesses first keep the call free. Irenaeus: God ``does indeed, through His infinite kindness, call the
unworthy; but He examines those who are called'' (IV.36.6). Hilary: the
calling ought to have made them good, \latin{quia sancta est}, because it is
holy, but \latin{per vitium inemendatae voluntatis discrimen est vocatorum},
through the fault of an unamended will there is a difference among the
called; the fault was one \latin{Deus solus invenit}, God alone finds (22.6--7).
Chrysostom: ``although to be called and to be cleansed was of grace, yet, when
called and clothed in clean garments, to continue keeping them so, this is of
the diligence of them that are called'' (69.2). Aquinas: \latin{noluit quod
mali venirent nisi pararent se}, he did not will that the bad should come
unless they prepared themselves (p.~282).

What the garment is, they say differently. Augustine searches by elimination.
It is not baptism, ``for this garment I see in the good, I see in the bad'';
nor the altar, nor fasting, nor miracles. It is ``charity out of a pure heart,
and of a good conscience, and of faith unfeigned'' (1~Tim 1:5): ``Let charity
be born in thee, if it be yet unborn, and if it be born, be it nourished,
fostered, increased'' (\work{Sermo} 90.5--6). And it is given: ``Clothe others,
and ye are clothed yourselves. \ldots\ Christ is naked; and He will give you
that `wedding garment'\,'' (\work{Sermo} 95.7). Gregory reaches the same answer
by the same route. If the garment were baptism or faith, who could have
entered without them? \latin{Quid ergo debemus intelligere nuptialem vestem,
nisi charitatem?} What then are we to understand by the wedding garment,
except charity? The believer lacks it \latin{si charitatis gratiam non
custodit}, if he does not keep the grace of charity (38.9).

The others name it in their own terms: for Chrysostom ``life and practice'';
for Irenaeus works of righteousness, ``so that the Spirit of God may rest upon
us''; for Jerome the Lord's commandments and the works that make the new man's
clothing (cols.~160--161). Aquinas gathers these into one: \latin{Quae est
ista vestis? Christus}, what is this garment? Christ. Some put him on by the
sacrament, some by charity, others by remembering death or by conforming their
works to his; and to have the wedding garment is to put him on \latin{per
operationem bonam, per conversationem sanctam, per caritatem veram: et si unum
deficiat, malum}, by good works, by holy living, by true charity, and if one
of these is lacking, it is bad (p.~282). Hilary's answer stands against Augustine's and
Gregory's: the garment is \latin{gloria Spiritus sancti}, the glory of the
Holy Spirit, received \latin{bonae interrogationis confessione}, at the
confession of the good questioning, and kept unstained to the kingdom (22.7).
He does not say baptism; the phrase recalls Peter's words on baptism (1~Pet
3:21), and on the editor's reading of that allusion his garment is received
at baptism and kept by the baptized. Augustine and Gregory say expressly that
it is not baptism, because the bad are baptized; Chrysostom's clean garment
given with the call stands between them, and Aquinas counts the sacrament
among the ways Christ is put on while the garment he requires is good works,
holy living and true charity together.

This disagreement is real, and the second interpretation does not settle it.
It claims what the witnesses share: the garment is required of those whom the
call has gathered, is to be kept, and is answered for when the king comes in;
and those who say where it comes from, Chrysostom, Augustine, Gregory and
Hilary, say that it is given, as the call is. Here the first interpretation's
two feasts return: Gregory's present Church of bad and good is the hall into
which the king comes, and Augustine says the closing saying shows ``who in
this present feast are accounted to be such, as to be brought to that other
feast, where no bad men shall come'' (90.4).

\subsection{``Friend'', the silence, and the few who are chosen}

In the longer form the king's question begins with a word the witnesses
weigh. Jerome: \latin{Amicum vocat, quod invitatus ad nuptias est}, he calls him
friend because he was invited (col.~161). Gregory: \latin{amice per fidem, sed
non amice per operationem}, friend by faith but not in deed; before that last
rebuke every argument of excuse falls silent (38.12). Yet Gregory tempers the
scene at once: whoever has the garment, but not yet perfectly, \latin{ad pii
regis ingressum desperare veniam non debet}, must not despair of pardon when
the merciful king comes in (38.12). That coming in is the judgment for Jerome,
Gregory and Aquinas, who adds death and the Church's tribulations (col.~160;
38.14; p.~282).

The witnesses come to the closing saying from different places, and not all
of them weigh how few the chosen are. Gregory passes to it from the guest cast
out, \latin{in quo videlicet omne malorum corpus exprimitur}, in whom the whole
body of the wicked is expressed (38.14), and had already told his hearers not
to be alarmed \latin{quod in Ecclesia et multi mali et pauci sunt boni}, with
the ark narrowing to one cubit, the strait way of Matthew 7:13--14 and the
threshing floor's few grains (38.8). Augustine states a proportion: the one
man ``not only was many, but those many in numbers far surpassed the number of
the good'' (90.4). Aquinas gives the saying its reason in the parable,
\latin{quia quidam nolunt venire, quidam non habent vestem nuptialem}, some will
not come and some lack the garment, and joins it to Matthew 7:14, ``few there
are that find it'' (p.~282). Irenaeus sets it beside the former covenant, in
which ``with many of them was He not well pleased'' (IV.36.6; 1~Cor 10:5).
Hilary puts the scarcity among the chosen, never in the call: \latin{Non est
igitur paucitas in invitatis, sed raritas in electis}, there is no fewness
among the invited, but a rarity among the chosen (22.7). Jerome reads the
saying otherwise: \latin{non initia, sed finis quaeratur}, not the beginning
but the end is sought (col.~161).

Hilary, Augustine, Irenaeus and Gregory keep the call itself without limit,
and Gregory turns the saying on his hearers: \latin{Quia vocati sumus, novimus; si sumus electi, nescimus}, that
we are called we know, whether we are chosen we do not know; so each must
humble himself, and \latin{omnes in sola divina misericordia gaudeant}, let all
rejoice in the divine mercy alone (38.14, 16). The saying itself names no
number; the witnesses' judgment of how few are chosen comes from the texts
they set beside it.

\subsection{To be full and to be hungry}

Chrysostom hears in Paul's ``I have learned'' a discipline ``not easy of
attainment'', and holds plenty as dangerous as want: ``Many know not how to be
full, as for example, the Israelites, `ate, and kicked'\,'' (Deut 32:15), and
the success is ``not mine own, but His who has given me strength''
(\work{Homilies on Philippians} 15). Aquinas makes the verses the
mark of the wise man, \latin{ille perfectus est qui scit uti quolibet statu},
he is perfect who knows how to use any state, because \latin{ex abundantia
erigitur animus contra Deum, ex paupertate deiicitur}, abundance raises the
mind against God and poverty casts it down (\work{Super Philippenses} c.~4,
lect.~2, pp.~401--403). No witness sets Paul beside the parable; heard beside
it in the third interpretation, he is the guest who would have come, not
without goods, for he has just received the Philippians' care, but not held by
them.

The reading's last verse serves two interpretations. Aquinas expounds the
Vulgate's \latin{impleat omne desiderium} first as God's return to the givers,
and then of glory, \latin{quia ibi implebitur totum desiderium}, because there
all desire will be filled, with the verse he gave Isaiah's third feast,
\latin{Satiabor cum apparuerit gloria tua}, ``I shall be satisfied when thy
glory shall appear'' (p.~403). The first interpretation hears in it the feast
completed, the third the end of its own line. Both rest on the Latin
\latin{desiderium}, which the Vulgate and the \work{Nova Vulgata} share and the
Douay--Rheims renders ``want''. Chrysostom's text is a blessing: Paul
``invokes blessings upon them, as poor men do''.

Where the first Communion antiphon is chosen, the third interpretation ends on
the rich who went hungry. Augustine meets the objection that ``many rich men
that are wicked die in their riches'' with the spiritual sense: ``Those rich
then lack, they lack, and what is heavier, they lack bread'', the bread of the
one who said, ``I am the Living Bread which came down from Heaven'', a joining
of the psalm to John that is his own (\work{Enarrationes in Psalmos} 33).
Robert Bellarmine gives the verse ``a higher meaning'': those attached to the
temporalities of this world ``always hunger and need'', while those who seek
the Lord, clinging ``to the Supreme Good'', possess all that is good
(\work{Commentary on the Book of Psalms}, Ps~33:10--11). The reading rests on
the Latin and Greek text; the Hebrew has young lions. Aquinas reads the
responsorial psalm's ``I shall want nothing'' of the future, \latin{quoniam
habebimus Deum}, because we shall have God (\work{Postilla} 22). In the first
interpretation the antiphon's rich lack the Living Bread; in the second they
stand against Augustine's heart ``filled with such jewelry of virtues''.

The Alleluia verse enters each line by a different comment. Chrysostom says the
hope of the calling is ``as yet \ldots\ hidden, but not so to the faithful'',
the condition of the invited who have not yet seen the feast, and that God
``would thus have us become good of our own will'' (\work{Homilies on
Ephesians} 3), the second interpretation's difficulty in a sentence. Aquinas
reads the enlightened eyes \latin{contra eos, qui habent oculos illuminatos
tantum ad temporalia cognoscenda}, against those whose eyes see only temporal
things (\work{Super Ephesios} c.~1, lect.~6), which is the third's.

\subsection{The Missal's prayers in three hearings}

No witness cited joins the Missal's antiphons and orations, which serve all
three years, to these readings; each interpretation hears them along its own
line. The Entrance Antiphon asks who could stand if the Lord marked
iniquities. Augustine asks, ``And what is this propitiation, except
sacrifice?'' (\work{Enarrationes in Psalmos} 129.3): in the first
interpretation, the door to the feast. Bellarmine says that without grace we
cannot even see ``the enormity of the offence'': in the second, why the guest
had nothing to say. Augustine says that those who prosper in their sins are
deepest in the depths, for ``a deceitful happiness is itself a greater
unhappiness'': in the third, the farm and the merchandise.

For the Collect's grace that goes before and follows, Aquinas founds the
division of grace on two psalm verses, the second this Sunday's ``thy mercy
will follow me'', and quotes Augustine: \latin{praevenit ut vocemur,
subsequitur ut glorificemur}, grace goes before that we may be called and
follows that we may be glorified (\work{Summa theologiae} I-II, q.~111, a.~3).
He does not cite the Collect. In the first interpretation the call is the
king's invitation and the glory the feast completed; in the second the grace
clothes the guest and keeps the garment on him, and Blessed Ildefonso
Schuster, commenting on the Gospel, holds call and garment together: ``God calls our souls, but they must be in
accord with their calling'' (\work{The Sacramentary}, vol.~3, p.~173); in the
third the prayer's \latin{intentos}, intent on good works, names the attention
the invited gave to farm and trade. The Prayer over the Offerings asks that
\latin{per haec piae devotionis officia} the faithful pass to heavenly glory:
the passage to the feast completed, the works of the garment offered, or goods
given rather than held.

Where the second Communion antiphon is chosen, the three hear it most
differently. In the first the guests see the King as he is, for Aquinas,
giving Origen's reason for the parable's ``man who is king'', says \latin{cum
videbitur sicuti est, tunc erit Rex}, when he is seen as he is, then he will
be King (p.~279), though he cites Paul's ``face to face'' for it. In the
second the vision is measured by charity: ``he will have a fuller
participation of the light of glory who has more charity'' (\work{Summa
theologiae} I, q.~12, a.~6), and Augustine, on the letter's next verse, which
the antiphon does not include, says ``Thou purifiest thyself, not by thyself,
but by Him who cometh to inhabit thee'' (\work{Homilies on the First Epistle of
John} 4.7). In the third it is desire stretched by waiting: ``if thou art full
of vinegar, where wilt thou put the honey?'' (4.6).

The Prayer after Communion asks that those fed with the Body and Blood of
Christ be made partakers of the divine nature. The fourth \work{Mystagogical
Catechesis}, transmitted under Cyril of Jerusalem's name, joins the two:
``we become partakers of the divine nature'' (4.3); and Aquinas calls grace
\latin{participatio divinae naturae} (I-II, q.~110, a.~3). In the first
interpretation the prayer asks that Gregory's restated wedding be shared at
the table; in the second the garment is the life of that grace; in the third
the verse of 2~Peter whose words it takes goes on, ``flying the corruption of
that concupiscence which is in the world'', a clause the prayer does not
quote.

\subsection{Where the three interpretations agree and part}

All three read the parable as a call wholly God's gift that can be refused or
wasted, and no witness cited denies another interpretation's claim; Gregory and Aquinas
contribute to all three from different sentences. They differ in the question put and so
in the verses that carry it: the first the \work{Ordo}'s correlation of
Isaiah with the Gospel, the second Matthew 22:11--14 with the Collect, the
third Matthew 22:5 with the second reading and the first Communion antiphon,
which neither the books nor any witness joins to the parable.

Their four senses part accordingly. Literally, the first reads the feast
promised and the hall filled, the second the guest inspected, the third the
excuses and Paul's contentment. Allegorically, the first finds the Son's
wedding with the Church; the second the garment, named charity, works, Christ
put on or the Spirit's glory; the third the mystery of the Incarnation refused
and the call passing to the gentiles at the outlets of the roads (Jerome) or to
those whom no prosperity held (Gregory). Morally, the first asks that the
present feast be received worthily, the second that charity be born and kept
in humility, the third that spiritual things be set above necessary ones.
Anagogically, the first ends in the wedding completed at the resurrection, the
second at the judgment and a vision measured by charity, the third in desire
filled in glory. Their consequences differ with them: hope and the universal
reach of the call, humility and the works of love, detachment and the right
use of goods.

Three real disputes cross these lines: what the garment is, dividing Augustine
and Gregory from Hilary, at the second interpretation's centre; whether the
wedding is the present Church or the glory of the resurrection, an emphasis
dividing Gregory and Hilary that shapes the first; and whether the burnt city
is Jerusalem under Titus or the eternal fire, dividing Chrysostom from Gregory
and Hilary, with Jerome and Aquinas giving both. Where the shorter Gospel is
read, the third interpretation holds whole, the first keeps its two feasts
joined by the resurrection alone, and the second keeps only the gathering of
bad and good. Each hears both Communion antiphons, but the first antiphon
weighs most in the third interpretation and the second in the first.
